\section{Conclusion}
\label{sec:conclusion}

We set out to answer a narrow question -- can a driving system detect the onset of snow-induced perception degradation with a false-alarm rate it can actually budget for, continuously, without a fixed monitoring horizon -- by wrapping a multimodal 3D detector in split conformal prediction and tracking its miscoverage rate with an anytime-valid e-process. That backbone is established prior work~\cite{MonroyMunoz2026WACV,WaudbySmithRamdas2024JRSSB}, not ours; what we attempted to add was a betting rule conditioned on the detector's own cross-modal disagreement signal as a leading indicator, a spatially resolved grid of per-cell e-processes with online false-discovery-rate control~\cite{WangRamdas2022JRSSB} in place of one global alarm bit, and an evaluation against Snowy Scenes~\cite{SnowyScenes2026}, a real adverse-weather dataset, rather than only synthetic corruption of a clear-weather one.

The headline real-data outcome is a negative one, and we report it as such rather than reframing it after the fact. On the corrected checkpoint, the covariate-informed bettor never alarms within the 20-frame window at any tested $\delta$, across all 5 onset replicates -- strictly worse than the covariate-blind baseline, which alarms at delay 3 ($\delta=0.3$) and delay 10 ($\delta=0.1$), unchanged from before the fix since it never consumes the CCP score at all. Section~\ref{sec:discussion} traces this to the corrected covariate varying in the wrong direction for this task (higher disagreement in the nominal \texttt{accumulated} regime than in the degraded \texttt{falling} regime) and to a plausible mechanism: PGD-adversarial perturbations, the mismatch signal the fix uses to teach $S$ to disagree, are not the same kind of feature-space disturbance real snowfall produces, and there is no reason a priori to expect a covariate trained against the former to transfer as a leading indicator of the latter. We do not know, from this evaluation alone, whether that transfer failure is fundamental to using an adversarially-supervised consistency signal for a naturally-occurring weather shift, or an artifact of an untuned $\kappa$, a small training run, or this specific detector architecture -- and we say so rather than picking whichever explanation makes the paper's story cleaner.

Getting real access to Snowy Scenes changed two parts of the plan we had not anticipated needing to change, independent of the result above. The dataset contains no clear-weather split at all, so our onset stream splices held-out mild-snow (\texttt{accumulated}) frames with heavy-snowfall (\texttt{falling}) frames rather than a genuine clear-to-snow transition -- a real, physically grounded severity contrast, but a constructed splice point, not a naturally occurring in-dataset event. We say this plainly rather than implying a stronger design than the data supports, and we note the alternative we did not have the means to try: splicing in real KITTI clear-weather frames as the nominal side, which trades a within-dataset class-space guarantee for a genuinely clear baseline. And training the detector on real data surfaced the CCP-collapse finding described in Section~\ref{sec:training} before we could even measure the result above -- the consistency score our covariate-informed bettor depends on had, before we fixed the training objective, no real signal in it at all. We think both findings -- the collapse, and the fixed covariate's wrong-direction transfer -- are worth a reader's attention independent of each other and independent of whether either specific diagnosis turns out to be the last word: reusing an internal model signal trained for one objective as an external statistical covariate for a different one is not automatically safe even once you have confirmed the signal varies, and checking not just \emph{whether} a covariate varies but \emph{whether it varies in the direction your method needs} is a second, easy-to-skip diagnostic that this paper's own experience argues should be standard practice.

We are equally direct about what this evaluation does not establish, beyond the central negative result above. It runs at a modest scale -- 5 onset replicates, 5 stationary-control replicates, 20-frame scenes, one detector architecture, one calibration split size, one fixed $\kappa$ -- and none of those choices are tuned; they are small-scale defaults carried over from earlier mock-data testing, not values selected to make the reported numbers look good, and Section~\ref{sec:protocol} is explicit that this is a pilot-scale, unpowered comparison rather than a statistically powered one. The spatial multiplicity-control grid (Section~\ref{sec:spatial}) is implemented, unit-tested end to end against a synthetic corruption stream, and confirmed correct on the same real detector and dataset used for the global operating curve -- but we have not yet run it against a genuinely multi-object real driving scene with real spatial heterogeneity in where the perception stack degrades, which is the setting the spatial construction is actually motivated by and the natural next experiment, and which the negative result above does not itself bear on (the spatial grid does not depend on the CCP covariate being informative unless a covariate-informed bettor is plugged into it). The onset-transition framing itself, as discussed above, is a scoping compromise a reviewer may reasonably push back on. And the competitive-timing risk we flagged going in remains real: a co-author of~\cite{MonroyMunoz2026WACV} has an existing research line in conformal prediction for multi-object 3D detection uncertainty in autonomous driving and is well positioned to make an AV/3D extension of their own work independently of ours.

Future work follows directly from those gaps: determining whether a covariate computed differently -- a discrepancy between the two modalities' own independently predicted object locations, say, rather than a jointly-trained consistency head supervised against adversarial perturbations -- transfers better to a naturally-occurring weather shift than the mechanism we tried; running the spatial e-BH grid against real multi-object Snowy Scenes frames rather than only the global monitor; tuning $\kappa$ and the calibration split size against real data rather than carrying over untuned defaults, if a future covariate does show a real signal in the right direction to tune against; scaling the replicate count and scene length well beyond what this first real evaluation used; and, if reviewers judge the accumulated-to-falling splice an insufficiently strong stand-in for a naturally occurring onset, attempting the cross-dataset KITTI-clear-to-Snowy-Scenes-snow splice as an alternative. We think the real-data evaluation protocol, the spatial monitoring construction, and the two documented failure modes -- the CCP collapse and its fix, and the fixed covariate's wrong-direction transfer -- are each a real contribution on their own, independent of the fact that neither the original mechanism nor its fix produced the positive result we set out to demonstrate; we have tried throughout this paper to be as clear about where the evidence for each claim currently stops as we are about what it currently shows.
